\worksheettitle
  {Названия и диагонали}
  {\modulemark{M06} Учебный модуль \textbullet\ полный маршрут}

\studentnameblock

\begin{notebox}[Чему вы научитесь]
  Проверять название многоугольника по обходу границы, различать соседние и
  несоседние вершины, стороны и диагонали.
\end{notebox}

\begin{checkpointbox}[Вход: проверьте опору]
  У четырёхугольника $ABCD$ стороны $AB$, $BC$, $CD$, $DA$.
  \begin{enumerate}[label=\textbf{\arabic*.}]
    \item Какая вершина соседняя с $A$: $B$ или $C$? \answerblank[15mm]
    \item Можно ли назвать фигуру $ACBD$? \answerblank[20mm]
    \item Как это проверить? \answerblank[65mm]
  \end{enumerate}
\end{checkpointbox}

\begin{notebox}[Как устроен маршрут]
  Сначала разберите правило обхода границы. Затем научитесь проверять
  соседство вершин, отличать сторону от диагонали, выполнять построения и
  называть части многоугольника. После ошибки пройдите контроль.
\end{notebox}

\section{Название задаёт обход}

\correctnamingfigure

\begin{fillbox}[Разберите пример]
  Названия $ABCDEF$, $BCDEFA$ и $AFEDCB$ верны: каждая следующая вершина
  соседняя с предыдущей. Название $ACDEFB$ неверно, потому что переход
  $A-C$ идёт не по стороне.
\end{fillbox}

\begin{notebox}[Правило проверки]
  Поставьте палец в первой названной вершине и двигайтесь по сторонам.
  Каждая следующая буква должна обозначать соседнюю вершину. Обход можно
  начинать с любой вершины и выполнять в двух направлениях.
\end{notebox}

\section{Практика с опорой}

\begin{practicebox}[1. Проверьте названия]
  Для шестиугольника на рисунке отметьте верные записи:
  \[
    \square\ CDEFAB\qquad
    \square\ FEDCBA\qquad
    \square\ ABDEFC\qquad
    \square\ DEFCBA.
  \]
  Для неверной записи назовите первый нарушенный переход:
  \answerblank[55mm].
\end{practicebox}

\section{Сторона или диагональ}

\diagonalfigure

\begin{fillbox}[Что сделали]
  Сторона соединяет \answerblank[30mm] вершины, а диагональ ---
  \answerblank[30mm] вершины многоугольника.

  Из вершины $A$ на рисунке проведены диагонали
  \answerblank[55mm].
\end{fillbox}

\begin{notebox}[Алгоритм]
  \textbf{1.} Найдите обе вершины отрезка.
  \textbf{2.} Проверьте, идут ли они подряд по границе.
  \textbf{3.} Соседние вершины соединяет сторона, несоседние --- диагональ.
\end{notebox}

\begin{practicebox}[2. Диагонали из одной вершины]
  Проведите все диагонали шестиугольника, выходящие из вершины $A$.
  \hexagondiagonalsstudent
  Запишите их: \answerblank[55mm]. Количество: \answerblank[15mm].
\end{practicebox}

\section{Самостоятельная практика}

\begin{practicebox}[3. Классифицируйте отрезки]
  В шестиугольнике $ABCDEF$ определите вид каждого отрезка:
  \[
    AB\text{ --- }\answerblank[22mm],\qquad
    AC\text{ --- }\answerblank[22mm],
  \]
  \[
    FA\text{ --- }\answerblank[22mm],\qquad
    CE\text{ --- }\answerblank[22mm].
  \]
\end{practicebox}

\pentagonsplitfigure

\begin{practicebox}[4. Назовите полученные части]
  Диагональ $AC$ разделила пятиугольник $ABCDE$ на два многоугольника.
  Назовите их, обходя каждую границу последовательно:
  \answerblank[45mm] и \answerblank[55mm].
\end{practicebox}

\section{Разбираем ошибку}

\begin{warningbox}[«Любой отрезок между вершинами --- диагональ»]
  Ученик назвал сторону $BC$ диагональю шестиугольника.

  \textbf{Причина ошибки:} \answerblank[100mm]

  \textbf{Исправление:} \answerblank[100mm]
\end{warningbox}

\section{Контрольная точка}

\hexagondiagonalsstudent

\begin{checkpointbox}[5. Выполните без подсказки]
  Для шестиугольника $ABCDEF$:
  \begin{enumerate}[label=\textbf{\arabic*.}]
    \item отметьте все правильные названия:
      $ABCDEF$, $CDEFAB$, $AEDCBF$, $AFEDCB$;
    \item проведите все диагонали из вершины $C$;
    \item запишите их обозначения; \answerblank[60mm]
    \item объясните, почему $CD$ не входит в этот список. \writinglines{2}
  \end{enumerate}
\end{checkpointbox}

\begin{reflectionbox}[Проверяю себя]
  \begin{itemize}[label=\choicebox]
    \item проверяю название обходом границы;
    \item различаю соседние и несоседние вершины;
    \item отличаю сторону от диагонали;
    \item называю части после проведения диагонали.
  \end{itemize}
\end{reflectionbox}

\begin{resultbox}[Дальнейший маршрут]
  Если нарушаете обход или принимаете сторону за диагональ, выполните M06-R.
  Для тренировки используйте M06-H. После устойчивого результата продолжайте
  следующий раздел. M06\(\star\) выполняется дополнительно.
\end{resultbox}
